\documentclass[10pt,a4paper]{amsart}
\usepackage[margin=19mm]{geometry}
\usepackage{amsmath,amssymb,mathtools}
\usepackage[scr=rsfs,cal=euler]{mathalfa}
\usepackage{xfrac}
\usepackage[unicode,hidelinks,bookmarks=false]{hyperref}
\newcommand{\ie}{i.e.\ }
\newcommand{\cf}{cf.\ }
\newcommand{\resp}{respectively\ }
\newcommand{\Subsection}{\subsection}
\providecommand{\nobreakdash}{\nobreak\hbox{-}}
\setlength{\emergencystretch}{2em}
\multlinegap0pt
\usepackage{bm}
\usepackage{bbm}
\usepackage{dsfont}
\usepackage{xspace}
\usepackage{cancel}
\usepackage{mathtools}
\usepackage{amsthm}
\makeatletter
\@ifundefined{theorem}{\newtheorem{theorem}{Theorem}}{}
\@ifundefined{proposition}{\newtheorem{proposition}[theorem]{Proposition}}{}
\makeatother
\title[Minimal value set binomials and Frobenius nonclassical curve]{Minimal value set binomials and Frobenius nonclassical curves}
\author{Tiago Aprigio, Jo\~ao Paulo Guardieiro}
\date{Source: arXiv:2508.16541v2 (CC BY 4.0)}
\begin{document}
\maketitle

\noindent\emph{Source and licence.} Minimal value set binomials and Frobenius nonclassical curves. Version arXiv:2508.16541v2 (CC BY 4.0), distributed under the Creative Commons Attribution 4.0 International licence, as recorded in the arXiv licence metadata for this exact version and shown on the abstract page https://arxiv.org/abs/2508.16541v2. Full rights notice: licenses/arxiv-2508.16541-CC-BY-4.0.txt.\par\smallskip

\noindent\textit{Editorial scope.} Counted unit: the Proposition whose source label is \texttt{Vf=2} in the article (source0.tex lines 273--279, source label \texttt{Vf=2}), delivered together with the article's own complete original proof of it (source0.tex lines 281--365, one \texttt{proof} environment of 85 lines). No mathematics is written, repaired or re-derived: statement and proof are the article's own text line for line. Required context retained uncounted: eq:binomial-equation (lines 249--251). Every definition, display and label of those lines is retained unchanged. Omitted from this package and not redistributed: 1--248, 252--272, 280--280, 366--1028. Nothing inside a retained excerpt is omitted or altered: every retained body is delivered exactly as the article's lines are written, so no omission receipt of any kind accompanies this package. Citations: the retained excerpts carry no citation command at all, so no bibliography entry of the article is transcribed and no reference is redistributed. Reference adaptation: none was needed and none was made; every reference inside the delivered unit resolves inside it, so no unresolved reference or citation is produced. Printed slips kept literal: no printed word, symbol, hypothesis or number of the counted statement or of its proof was altered. Layout and class: the article's own class is \texttt{elsarticle} with packages including inputenc, amssymb, amsmath, amsthm, xcolor, enumitem, tikz, hyperref, ulem; the wrapper is the lane's pinned \texttt{amsart} editorial wrapper at 10pt on A4, which restates the theorem-like environments the retained text needs and adds the article's own macro definitions that the retained text needs (none), with extra wrapper packages (bm, bbm, dsfont, xspace, cancel, mathtools). The delivered numbering is the unit's own. Assets: no retained span contains a figure environment, a graphics-inclusion command, a PDF-page inclusion, a table or a TikZ picture; no image file is copied into this package, the package declares no asset dependency and no image file is redistributed with it. Licence: version arXiv:2508.16541v2 of this article is distributed under the Creative Commons Attribution 4.0 International licence, as recorded in the arXiv licence metadata for this exact version and shown on the abstract page https://arxiv.org/abs/2508.16541v2. A full rights notice is delivered at licenses/arxiv-2508.16541-CC-BY-4.0.txt. Characters: every retained excerpt is pure ASCII, so the delivered engine, XeLaTeX with its default Latin Modern text font, reports no missing character.
\begin{equation}\label{eq:binomial-equation}
F(x) = x^a + \beta x^b, \quad a > b, \quad \beta \in \mathbb{F}_q^*.
\end{equation}

\begin{proposition}\label{Vf=2}
    Let $F(x)$ be a binomial over $\mathbb{F}_q$. Then $F$ is an MVSP with $\# V_F = 2$ if, and only if, one of the following holds:
    \begin{enumerate}
        \item[$\mathrm{(i)}$] $F(x) = x^{q-1} + \beta x^{\frac{q-1}{2}}$, where $p > 2$ and $\beta^2 =1$. In this case, we have $V_F = \{0, 2\}$.
        \item[$\mathrm{(ii)}$] $F(x) = x^{\frac{2(q-1)}{3}} + \beta x^{\frac{q-1}{3}}$, where $p = 2$, $q > 2$ and $\beta^3 = 1$. In this case, we have $V_F = \{0, \beta^2\}$.
    \end{enumerate}
\end{proposition}

\begin{proof}
The verification that the polynomials in {\rm (i)} and {\rm (ii)} are MVSPs with the stated value set is straightforward.

Conversely, assume that $F(x) = x^a + \beta x^b$ is an MVSP over $\mathbb{F}_q$ with $\#V_F = 2$. From the definition of MVSPs, we have
\[
\frac{q-1}{2} < a \leq q-1,
\]
which implies $q > 2$. For $q = 3$, a direct computation shows that the only MVSP is $x^2 + \beta x$, with $\beta = \pm 1$.
For $q = 4$, the only MVSPs are of the form $x^2 + \beta x$, with $\beta \in \mathbb{F}_4^*$. Henceforth, we assume $q \ge 5$.

Let $\xi$ be a primitive element of $\mathbb{F}_q$. For all $k \ge 2$, the sequence $\{F(\xi^k)\}$ satisfies the recurrence relation
\begin{equation}\label{eqxi}
F(\xi^k) = (\xi^a + \xi^b)F(\xi^{k-1}) - \xi^{a+b}F(\xi^{k-2}).
\end{equation}

\medskip

\noindent
\textbf{(i): $p \neq 2$.} We need to consider some sub cases.

\noindent
\textbf{(i.1): The case $F(\xi) = 0$.}
In this case, $\xi^{a-b} = -\beta$. Since $a-b < q-1$ and $\xi$ is a primitive element, we have $\beta \neq -1$, and therefore $V_F = \{0, 1+\beta\}$. From Equation \eqref{eqxi} for $k=2$, we obtain $F(\xi^2)=1+\beta$ and $\xi^{a+b}=-1$. Thus,
\[
F(\xi^k) = (\xi^a+\xi^b)F(\xi^{k-1}) + F(\xi^{k-2}).
\]
In particular, $F(\xi^3) = (\xi^a + \xi^b)F(\xi^2)$ and we have $\xi^a + \xi^b$ equals $0$ or $1$.

If $\xi^a+\xi^b = 1$, then $F(\xi^3) = F(\xi^2)$, so $F(\xi^4)=2F(\xi^2)\notin V_F$, a contradiction. Hence, $\xi^a+\xi^b=0$, and this implies
\[
F(\xi^{2n}) = F(\xi^2), \qquad F(\xi^{2n+1}) = F(\xi) = 0,
\]
for all $n \ge 1$. Since \( x^{a - b} + \beta \) has exactly \( \gcd(a - b, q - 1) \) distinct roots in \( \mathbb{F}_q \), and \( a \leq q - 1 \), we can conclude that
$$
a = b + \frac{q - 1}{2}.
$$
From $F(\xi^2)=\xi^{2b}(1+\beta)$, we deduce $b = \frac{q-1}{2}$ and $\beta = 1$.

\smallskip

\noindent
\textbf{(i.2): The case $F(\xi) \neq 0$.}
If $\beta \neq -1$, then $F(\xi)=1+\beta$ and
\[
F(\xi^2) = (\xi^a+\xi^b-\xi^{a+b})F(\xi).
\]
Therefore, $\xi^a+\xi^b-\xi^{a+b}$ equals $0$ or $1$. In the first case, we obtain $F(\xi^2) = 0$, $F(\xi^3) = -\xi^{a + b} F(\xi)$ (so $\xi^{a + b} = -1$) and $F(\xi^4) = -F(\xi^3) \notin \{0, F(\xi)\}$, which is a contradiction. In the second case, we will obtain $F(\xi^k) = F(\xi) = 1 + \beta$ for all $k \geq 1$. Hence, the polynomial $F(x) - (1 + \beta)$ has at least $q - 1$ distinct roots in $\mathbb{F}_q$, which implies $a = q - 1$ and
  $$
  F(x) - (1 + \beta) = \prod_{i = 1}^{q - 1}(x - \xi^i) = x^{q - 1} - 1,
  $$
  which contradicts Equation \eqref{eq:binomial-equation}. Therefore, this case also does not occur. 
  
  So, if $F(\xi) \neq 0$, we must have $\beta = -1$. In this situation, necessarily $\xi^a+\xi^b=0$ (otherwise we would obtain $F(\xi^4) = - F(\xi) \notin V_f$), which yields, from Equation \eqref{eqxi}
\[
F(\xi^{2n})=0, \qquad F(\xi^{2n+1})=F(\xi),
\]
for all $n\ge1$. Arguing as before, we obtain
\[
a = b + \frac{q-1}{2}, \qquad b = \frac{q-1}{2}.
\]

\medskip

\noindent
\textbf{Case (ii): $p = 2$.} The proof is analogous to case (i). Therefore, we will omit some details.

\noindent
\textbf{(ii.1): The case $F(\xi) = 0$.} We conclude again that $\beta \neq 1$ and $\xi^{a+b} = 1$ (in particular, $a + b = q - 1$). Since the characteristic of the field is $2$, we have $\xi^a + \xi^b \neq 0$. From $F(\xi^3) = (\xi^a + \xi^b)(1 + \beta)$, we obtain $\xi^a + \xi^b = 1$ and
$$F(\xi^k) = F(\xi^{k - 1}) + F(\xi^{k - 2})$$
for all $k \geq 2$. This implies:
$$F(\xi^{3n + 1}) = 0 \quad \text{and} \quad F(\xi^{3n}) = 1 + \beta = F(\xi^{3n+2}) \quad \text{for all } n \geq 0.$$
Hence, the polynomial $F(x)$ now have $\frac{q-1}{3}$ roots in $\mathbb{F}_q$. We know that the number of its roots is $\gcd(a - b, q - 1)$ and that $a + b = q - 1$. We then obtain
$$a = \frac{2(q - 1)}{3}, \quad b = \frac{q - 1}{3}$$
and so $\beta = \xi^{a - b} = \xi^{(q-1)/3}$

\smallskip

\noindent
\textbf{(ii.2): The case $F(\xi) \neq 0$.} If $\beta = 1$, we analyze $F(\xi^2)$ and $F(\xi^3)$ to conclude that $\xi^a + \xi^b = 1$, $\xi^{a+b} = 1$ and $\xi^{3(a - b)} = 1$. We obtain again 
$$(a, b) = \left( \frac{2(q - 1)}{3},\, \frac{q - 1}{3} \right) \quad \textrm{and} \quad  \beta^3 = 1.$$

If $\beta \neq 1$, then $F(\xi^2) = (\xi^a + \xi^b + \xi^{a+b})F(\xi)$, and so $\xi^a + \xi^b + \xi^{a+b}$ equals $0$ or $1$. In the first case, we analyze $F(\xi^3)$, $F(\xi^4)$ and $F(\xi^5)$ to conclude that $\xi^{a+b} = \xi^{3(a - b)} = 1$. As before, we obtain
$$(a, b) = \left( \frac{2(q - 1)}{3},\, \frac{q - 1}{3} \right) \quad \textrm{and} \quad  \beta^3 = 1.$$
In the second case, we will obtain $F(\xi^k) = F(\xi)$ for all $k \geq 1$. Analogously to one of the sub cases of (i.2), we get a contradiction with Equation \eqref{eq:binomial-equation}.
\end{proof}

\end{document}
